\section{Conclusiones y Trabajos Futuros}

\subsection{8.1 Conclusiones}

Particularmente llamativo resulta el cierre de un ciclo de investigación que comenzó con desafíos bien conocidos y culmina con un prototipo que supera varias expectativas iniciales. Este trabajo demostró que la combinación de una topología CLCLC resonante optimizada con control predictivo por modelo y modulación Triple Phase-Shift representa una vía efectiva para elevar significativamente el rendimiento de los sistemas de potencia eléctrica.

Hassan et al. \cite{Hassan2026} establecieron un referente sólido con convertidores CLCLC experimentales de alta eficiencia. La presente propuesta iguala esos niveles (alcanzando picos de 97.8\%) mientras mejora de forma notable la respuesta dinámica y la integración sistémica. Se validó experimentalmente en un prototipo de 1 kW que mantiene conmutación suave en amplio rango, excelente estabilidad y eficiencia superior al 96\% en operación bidireccional.

Los resultados confirman que el modelado en estado-espacio detallado, el diseño cuidadoso de resonantes y la implementación digital híbrida DSP+FPGA constituyen una metodología reproducible. Más allá de los números, el sistema exhibe una inteligencia operativa que lo hace adecuado para microgrids DC con alta penetración de renovables.

\subsection{8.2 Limitaciones}

Lejos de ser un asunto resuelto, toda solución técnica trae consigo compromisos que merecen reconocimiento explícito. Reddy et al. \cite{Reddy2025} identifican con claridad los desafíos persistentes en integración de convertidores para movilidad y renovables; varios de ellos se manifestaron durante este desarrollo.

La complejidad computacional del MPC limitó el horizonte de predicción, obligando a balancear desempeño y viabilidad en tiempo real. Aunque se logró excelente comportamiento dinámico, a potencias muy bajas aparece degradación parcial de la conmutación suave. Además, el prototipo de 1 kW, aunque representativo, deja abierta la cuestión de escalabilidad a decenas de kilowatts donde efectos térmicos y parasíticos se acentúan.

No obstante, cabe matizar que estas limitaciones no invalidan los hallazgos, sino que delimitan con precisión el alcance actual de la solución. Factores externos como variaciones extremas de temperatura o envejecimiento de componentes no fueron evaluados exhaustivamente, aspecto común en trabajos académicos pero relevante para aplicaciones industriales.

\subsection{8.3 Trabajos Futuros}

Uno de los desafíos persistentes que surge al concluir proyectos de este tipo es resistir la tentación de declarar victoria prematura. Queda mucho terreno por explorar.

Mezouari et al. \cite{Mezouari2025} destacan en su revisión la necesidad de avanzar hacia topologías más densas y controles adaptativos. En esa dirección, se propone extender el esquema actual incorporando control adaptativo o MPC con aprendizaje por refuerzo para manejar variaciones paramétricas en tiempo real. La escalabilidad a potencias superiores (5-50 kW) con dispositivos SiC/GaN de nueva generación constituye otra línea prioritaria.

Resulta prometedor investigar la integración con técnicas de diagnóstico en línea de fallas y modos de operación resilientes ante contingencias de red. Asimismo, la optimización multiobjetivo en tiempo real a nivel de microgrid completa, incluyendo economía de operación y vida útil de baterías, abriría aplicaciones más amplias en sistemas de distribución moderna.

Reddy et al. \cite{Reddy2025} sugieren que el futuro de los convertidores DC-DC pasa por mayor inteligencia distribuida y densidad de potencia. La presente investigación proporciona una base sólida desde la cual abordar estos retos con mayor confianza. Los próximos pasos ya están definidos: prototipos de mayor potencia, pruebas de campo y exploración de algoritmos de control híbridos que combinen lo mejor de la predictividad y la adaptabilidad.